\chapter{数学基础；逻辑；集合论。}

对通常所说的“数学基础”的研究，自 19 世纪初以来一直在不间断地进行，但若没有一个使逻辑系统化的平行努力——至少是使逻辑中支配数学命题之间联系的那些部分系统化——这项研究是不可能取得成果的。因此，不可能把集合论的历史和数学形式化的历史同“数理逻辑”的历史分开来讲述。然而，传统逻辑如同近代哲学家的逻辑一样，其原则上的适用范围远远超出数学。所以读者不要指望在下文中找到一部逻辑学史，哪怕是以极其简略的形式写成的；我们只限于尽可能追溯逻辑的演变，而且仅限于它触及数学演变的那些部分。正是由于这个原因，对于非经典逻辑（多值逻辑、模态逻辑），我们将只字不提；我们更无法处理那些论争的历史，这些论争从智者派一直延续到维也纳学派，从未停止使哲学家们发生分裂——分歧既在于把逻辑应用于现实世界的对象或人类思想的概念是否可能，也在于如何应用。

在希腊时代之前就存在一种相当发达的数学，这在今天已无任何疑问。不仅整数和量的度量这些（已经非常抽象的）概念，在我们所掌握的来自埃及或迦勒底的最古老的文献中已被普遍使用，而且巴比伦代数凭借其方法的优雅与可靠，也不应被看作一堆靠经验摸索解出的问题的简单汇集。再者，虽然在这些文本中找不到任何与该词形式意义相近的“证明”，但有理由相信，这类解法——其一般性是通过特殊的数值应用显现出来的——的发现，若无最低限度的逻辑联系是不可能的（这些联系也许并非完全自觉，而颇像现代代数学家在着手一项计算时、在“形式地写下”全部细节之前所依赖的那些联系）（{[}232{]}，第~203~页及以下）。

希腊人本质的独创性，恰恰在于一种自觉的努力：把数学证明排成这样一个序列，使得从一个环节过渡到下一个环节不留下任何怀疑的余地，并迫使人们普遍同意。

希腊数学家在研究过程中，正如现代数学家一样，使用的是“启发式”的而非令人信服的论证，这一点，只要看看阿基米德的《方法》一文 {[}153 c{]}便可以证明（如果还需要证明的话）；其中还值得注意的是，它提到了早期数学家“已发现但未证明”的结果。\(^{1}\) 不过，从我们所知的最早一批详细文本（其年代为公元前 5 世纪中叶）起，数学文本的理想“典范”便已完全确立。它将在伟大的经典作家欧几里得、阿基米德和阿波罗尼奥斯那里达到最高的表现；在这些作者那里，证明的概念与我们的证明概念毫无二致。

我们没有任何文本可以据以追溯这种“演绎方法”的最初步骤；当我们刚一觉察到它的存在时，它在我们看来已近乎完善。人们只能认为，它相当自然地契合于希腊思想所特有的、对世界“解释”的永不停歇的追求，而这种追求早在公元前 7 世纪的爱奥尼亚哲学家那里便清晰可辨；此外，传统一致地把这种方法的发展和完善归功于毕达哥拉斯学派，其年代介于公元前 6 世纪末与公元前 5 世纪中叶之间。

正是这种对其目标和方法都有充分自觉的“演绎”数学，成为后世哲学思想和数学思想集中之所在。我们将看到，一方面，以数学为蓝本的“形式”逻辑逐步建立起来，并最终以形式化语言的创立而告终；另一方面，主要从 19 世纪初开始，数学的基本概念受到越来越多的质疑，人们将为澄清其本性作出巨大努力，尤其是在集合论诞生之后。

